\section{Logros y umbrales del aprendizaje por refuerzo}

\subsection{Filosofía y método de la serie Zero}
Hubert atribuye el éxito de AlphaZero a dos ejes centrales: \textbf{search-integrated learning} y \textbf{reliable reward from simulation}. En los juegos, el entorno simulado ofrece perfect feedback, de modo que el RL ya no depende de etiquetas humanas, sino que usa las victorias y derrotas de la búsqueda como verdad de referencia.
\begin{knowledgebox}{Los tres pilares de la filosofía Zero}
1) Self-play generating curriculum; 2) Search results distilled into policy/value; 3) Preference for small, reusable modules rather than monolithic black-box policies.
\end{knowledgebox}

Él dijo: «El esplendor de la serie Zero no es el éxito de un solo modelo, sino de un proceso de ingeniería depurado con precisión: cada self-play produce un curriculum, y la policy/value se reconstruye como un heuristic module listo para la búsqueda». Esta filosofía influye directamente en las expectativas de policy distillation de AlphaProof.

\subsection{El volante de búsqueda-aprendizaje}
AlphaZero/AlphaTensor construyeron un volante de «exploración → aprendizaje → optimización». Hubert señala en particular: sin strong search no se obtienen high-quality trajectories; sin fast learning no es posible interiorizar el conocimiento nuevo; y sin una estrecha combinación entre policy y value, la search cae en un thousand-node dead-end.
\begin{warningbox}{Vacíos del volante y fallas potenciales}
1) Search de baja calidad: el agent no logra encontrar soluciones sólidas y la actualización de policy es débil; 2) reward con ruido: la etapa de aprendizaje incorpora hallucinated trajectories; 3) Compute limitado: una search depth insuficiente impide la convergence del learning.
\end{warningbox}

Hubert compara este volante con «excavar junto a un agujero negro»: la search no puede quedarse en la superficie, o el learning nunca tendrá con qué avanzar; el learning también debe absorb a tiempo los nuevos lemma que descubre la search, o la search seguirá repitiendo el mismo camino. El equipo de AlphaProof usa un two-tier policy buffer para que search-area y learning-rate se sigan mutuamente y evitar así el drift.

\subsection{Exploración, retroalimentación y generalización}
La retroalimentación de AlphaZero es un win-rate determinístico; AlphaProof, en cambio, debe enfrentar la «corrección formal de la demostración»: si la prueba en Lean no pasa, la retroalimentación es simplemente un fallo. Hubert explica: «En matemáticas, la reward es +1 o es 0, no hay zona gris, y eso obliga al RL a desarrollar trace-aware regret minimization». Comparación entre búsqueda y heurísticas:
\begin{table}[H]
\centering
\begin{tabular}{p{4cm}p{4cm}p{4cm}}
\toprule
Dimensión & Search-driven RL & Heuristic/Scripted \\
\midrule
Calidad de la retroalimentación & Verificación perfecta, replay buffer de alta precisión & Señal con mucho ruido, propensa al drift \\
Capacidad de generalización & Logra la transferencia entre teoremas mediante policy distillation & Requiere ingeniería manual \\
Capacidad de auditoría & Se puede reproducir la trajectory y ofrecer el proof term & Difícil de reducir a pasos verificables \\
\bottomrule
\end{tabular}
\caption{Comparación entre la retroalimentación de RL y los métodos heurísticos en AlphaProof}
\end{table}

\subsection{Resumen de la sección}
Los tres planos del éxito de AlphaZero (search, feedback, learning) son el punto de partida de AlphaProof; solo prestando atención a cualquier ruptura del volante, al colapso de la señal de retroalimentación y a los límites de la generalización se puede asegurar que el RL agent avance de forma estable en el mundo matemático.
\newpage

